% =========================================================
% Chains and Antichains
% =========================================================

\section{Chains and Antichains}
\label{sec:chains-antichains}

\begin{remark*}[Toolkit: Chains and Antichains]
A poset can contain subsets whose elements are either completely aligned
or entirely incomparable.  These two extremes are formalised as chains and
antichains.  They measure the ``linearity'' and ``width'' of the order
structure respectively.
\end{remark*}

% ---------------------------------------------------------
% def:lattice-chain
% ---------------------------------------------------------

\begin{definitionbox}{Definition (Chain)}
\begin{definition}[Chain]
\label{def:lattice-chain}
Let $(S, \leq)$ be a partially ordered set.  A subset $C \subseteq S$ is a
\emph{chain} if every pair of elements in $C$ is comparable:
\[
\forall x, y \in C, \quad x \leq y \;\lor\; y \leq x.
\]
Equivalently, the restriction of $\leq$ to $C$ is a
\hyperref[def:lattice-order-total-order]{total order} on $C$.
\end{definition}
\end{definitionbox}

\begin{remark*}[Standard quantified statement]
\[
\operatorname{Chain}(C, S, \leq)
\;\coloneqq\;
C \subseteq S
\;\land\;
\forall x, y \in C,\; x \leq y \lor y \leq x.
\]
\end{remark*}

\begin{remark*}[Negated quantified statement]
\[
\neg\operatorname{Chain}(C, S, \leq)
\;\coloneqq\;
\exists x, y \in C,\; x \not\leq y \land y \not\leq x.
\]
\end{remark*}

\begin{remark*}[Interpretation]
A chain is a subset in which the partial order behaves as a total order:
no two elements are left incomparable.  The entire real line $\mathbb{R}$
under its standard ordering is itself a chain.  Any finite totally ordered
set --- a ranking, a sequence of nested intervals, a linearly ordered list
of stages --- is a chain.  Chains model the ``ladder'' structure of order:
there is a clear direction from smaller to larger, and the path from any
element to any other is unambiguous.
\end{remark*}

\begin{dependencies}
\hyperref[def:lattice-order-partial-order]{Partial order},
\hyperref[def:lattice-order-total-order]{total order}.
\end{dependencies}

% ---------------------------------------------------------
% def:antichain
% ---------------------------------------------------------

\begin{definitionbox}{Definition (Antichain)}
\begin{definition}[Antichain]
\label{def:antichain}
Let $(S, \leq)$ be a partially ordered set.  A subset $A \subseteq S$ is
an \emph{antichain} if no two distinct elements of $A$ are comparable:
\[
\forall x, y \in A, \quad x \neq y \Longrightarrow \neg(x \leq y) \land \neg(y \leq x).
\]
\end{definition}
\end{definitionbox}

\begin{remark*}[Standard quantified statement]
\[
\operatorname{Antichain}(A, S, \leq)
\;\coloneqq\;
A \subseteq S
\;\land\;
\forall x, y \in A,\; x \neq y \Rightarrow x \not\leq y \land y \not\leq x.
\]
\end{remark*}

\begin{remark*}[Negated quantified statement]
\[
\neg\operatorname{Antichain}(A, S, \leq)
\;\coloneqq\;
\exists x, y \in A,\; x \neq y \land (x \leq y \lor y \leq x).
\]
\end{remark*}

\begin{remark*}[Interpretation]
An antichain is a subset in which the partial order makes no comparisons at
all: every pair of distinct elements is incomparable.  In the power set
$\mathcal{P}(X)$ ordered by inclusion, the collection of all subsets of
a fixed cardinality $k$ forms an antichain --- no subset of size $k$ contains
another.  Antichains measure the ``width'' of a poset: Dilworth's theorem
states that the minimum number of chains needed to cover a finite poset
equals the maximum size of an antichain.  In higher-dimensional structures
such as product orders, large antichains arise naturally and are the primary
tool for understanding the lateral spread of the order.
\end{remark*}

\begin{dependencies}
\hyperref[def:lattice-order-partial-order]{Partial order}.
\end{dependencies}
